Vývoj \ac{HDP} na obyvatele v~paritě kupní síly (\ac{PPS}, EU27\,=\,100) zasazuje českou ekonomiku do středoevropského konvergenčního příběhu: ČR se od roku 2000 posunula ze~zhruba \SI{70}{\percent} průměru EU27 na přibližně \SI{95}{\percent}, přičemž tempo růstu v~posledním desetiletí zpomalilo.
Tato relativně příznivá makroekonomická pozice však kontrastuje s~přetrvávajícími mzdovými zaostáváním zachyceným v~obr.~\ref{fig:net_income_ratio_timeline}: ekonomický výkon konverguje, ale čisté příjmy pracovníků nikoliv ve stejném tempu.
Srovnání s~AT a~DE --- zeměmi s~podobnou průmyslovou strukturou, ale s~rozvinutým kolektivním vyjednáváním na odvětvové úrovni --- naznačuje, že část tohoto „překladu” \ac{HDP} do mezd zprostředkovávají právě silné odbory a~rozšiřitelné kolektivní smlouvy (\ac{KSVS}).
Bez systémového mechanismu sdílení produktivitního přebytku (wage-indexation, sektorové \ac{KS}) zůstane mezera mezi výkonností a~odměňováním v~ČR strukturálním problémem i při pokračující nominální konvergenci.

Vývoj \ac{HDP} na obyvatele v~paritě kupní síly (\ac{PPS}, EU27\,=\,100) zasazuje českou ekonomiku do středoevropského konvergenčního příběhu: ČR se od roku 2000 posunula ze~zhruba \SI{70}{\percent} průměru EU27 na přibližně \SI{95}{\percent}, přičemž tempo růstu v~posledním desetiletí zpomalilo.
Tato relativně příznivá makroekonomická pozice však kontrastuje s~přetrvávajícími mzdovými zaostáváním zachyceným v~obr.~\ref{fig:net_income_ratio_timeline}: ekonomický výkon konverguje, ale čisté příjmy pracovníků nikoliv ve stejném tempu.
Srovnání s~AT a~DE --- zeměmi s~podobnou průmyslovou strukturou, ale s~rozvinutým kolektivním vyjednáváním na odvětvové úrovni --- naznačuje, že část tohoto „překladu” \ac{HDP} do mezd zprostředkovávají právě silné odbory a~rozšiřitelné kolektivní smlouvy (\ac{KSVS}).
Bez systémového mechanismu sdílení produktivitního přebytku (wage-indexation, sektorové \ac{KS}) zůstane mezera mezi výkonností a~odměňováním v~ČR strukturálním problémem i při pokračující nominální konvergenci.

Vývoj \ac{HDP} na obyvatele v~paritě kupní síly (\ac{PPS}, EU27\,=\,100) zasazuje českou ekonomiku do středoevropského konvergenčního příběhu: ČR se od roku 2000 posunula ze~zhruba \SI{70}{\percent} průměru EU27 na přibližně \SI{95}{\percent}, přičemž tempo růstu v~posledním desetiletí zpomalilo.
Tato relativně příznivá makroekonomická pozice však kontrastuje s~přetrvávajícími mzdovými zaostáváním zachyceným v~obr.~\ref{fig:net_income_ratio_timeline}: ekonomický výkon konverguje, ale čisté příjmy pracovníků nikoliv ve stejném tempu.
Srovnání s~AT a~DE --- zeměmi s~podobnou průmyslovou strukturou, ale s~rozvinutým kolektivním vyjednáváním na odvětvové úrovni --- naznačuje, že část tohoto „překladu” \ac{HDP} do mezd zprostředkovávají právě silné odbory a~rozšiřitelné kolektivní smlouvy (\ac{KSVS}).
Bez systémového mechanismu sdílení produktivitního přebytku (wage-indexation, sektorové \ac{KS}) zůstane mezera mezi výkonností a~odměňováním v~ČR strukturálním problémem i při pokračující nominální konvergenci.

Vývoj \ac{HDP} na obyvatele v~paritě kupní síly (\ac{PPS}, EU27\,=\,100) zasazuje českou ekonomiku do středoevropského konvergenčního příběhu: ČR se od roku 2000 posunula ze~zhruba \SI{70}{\percent} průměru EU27 na přibližně \SI{95}{\percent}, přičemž tempo růstu v~posledním desetiletí zpomalilo.
Tato relativně příznivá makroekonomická pozice však kontrastuje s~přetrvávajícími mzdovými zaostáváním zachyceným v~obr.~\ref{fig:net_income_ratio_timeline}: ekonomický výkon konverguje, ale čisté příjmy pracovníků nikoliv ve stejném tempu.
Srovnání s~AT a~DE --- zeměmi s~podobnou průmyslovou strukturou, ale s~rozvinutým kolektivním vyjednáváním na odvětvové úrovni --- naznačuje, že část tohoto „překladu” \ac{HDP} do mezd zprostředkovávají právě silné odbory a~rozšiřitelné kolektivní smlouvy (\ac{KSVS}).
Bez systémového mechanismu sdílení produktivitního přebytku (wage-indexation, sektorové \ac{KS}) zůstane mezera mezi výkonností a~odměňováním v~ČR strukturálním problémem i při pokračující nominální konvergenci.

\input{texparts/python/gdp_ppp_timeline}



